\section{Derivatives with respect to a scalar shift}
\label{sec:shifted_power_derivative}

Let $I$ be a finite index set and let $A\in\C^{I\times I}$ be positive
semidefinite. Norms are Euclidean operator norms. The matrix $A$ may be
singular, and $I$ may be empty.

% Independent supporting facts for OpenAI, Polynomial PEPS approximation
% of gapped square-grid ground states, September 24, 2026, 03-patches.tex,
% lines 425--451, at adc7f1241b42e322a6451854ab7e4b4c146bf78a.
% No assertion about the minimum envelope or regulator growth is made here.

\begin{theorem}[Derivative of a shifted power]
    \label{thm:shifted_power_derivative}
    \lean{Matrix.PosSemidef.hasDerivAt_add_smul_one_rpow}
    \leanok
    For every $b>0$ and $r\in\R$, the function $s\mapsto(A+s\Id)^r$
    has derivative $r(A+b\Id)^{r-1}$ at $b$.
\end{theorem}

\begin{proof}
    \leanok
    \uses{thm:shifted_power_functional_calculus}
    Fix a unitary eigenbasis $A=U\operatorname{diag}(\lambda_i)U^*$.
    For $s>0$, the shifted power is
    $U\operatorname{diag}((\lambda_i+s)^r)U^*$.
    Each scalar coordinate has derivative
    $r(\lambda_i+b)^{r-1}$ because $\lambda_i+b>0$.
    The maps taking diagonal entries to a matrix and multiplying by the
    fixed matrices $U$ and $U^*$ are continuous linear maps. Applying them
    to the coordinate derivatives gives the claimed derivative.
\end{proof}

\begin{theorem}[Exponential regularization]
    \label{thm:exp_regularization_derivative}
    \lean{Matrix.PosSemidef.hasDerivAt_add_exp_neg_smul_one_rpow}
    \leanok
    For every $R,r\in\R$, the function $s\mapsto(A+e^{-s}\Id)^r$
    has derivative $-re^{-R}(A+e^{-R}\Id)^{r-1}$ at $R$.
\end{theorem}

\begin{proof}
    \leanok
    \uses{thm:shifted_power_derivative}
    The chain rule and $\frac{d}{dR}e^{-R}=-e^{-R}$ give the assertion.
\end{proof}

\begin{theorem}[Shifted inverse contraction]
    \label{thm:shifted_inverse_contraction}
    \lean{Matrix.PosSemidef.l2_opNorm_smul_add_smul_one_rpow_neg_one_le}
    \leanok
    For every $b>0$, one has $\|b(A+b\Id)^{-1}\|\le1$.
\end{theorem}

\begin{proof}
    \leanok
    \uses{thm:shifted_power_functional_calculus}
    The matrix is the functional calculus of $t\mapsto b/(t+b)$ at $A$.
    Every eigenvalue $t$ of $A$ is nonnegative, so $0\le b/(t+b)\le1$.
    The operator norm bound follows from this scalar bound on the spectrum.
\end{proof}
